% !TEX root = ../Thesis.tex
\chapter{METHODOLOGY}
\section{Iterative development}
The project follows a prototype--test--improve approach. Initial learning content and interface elements are implemented first; subsequent iterations should respond to observed user difficulties and feedback. Changes will be documented in the project's Git repository.

\section{Implementation environment}
The prototype is built using React, JavaScript/JSX, CSS, and Vite. Collaboration and version history are managed through Git and GitHub. Code changes are intended to be checked with an automated production build.

\section{Evaluation design}
Planned evaluation includes task-based usability testing, basic accessibility checks, and knowledge questions administered before and after the learning activity where feasible. The team should record participant characteristics, number of participants, tasks, instruments, and the dates of testing before interpreting results.

\section{Ethical and methodological considerations}
Participation should be voluntary. Avoid collecting unnecessary personal data. Report limitations related to sample size, question design, and incomplete testing. This chapter describes planned methods; it is not a claim that participant testing has occurred.
